Die kinetische Energie $E_\mathrm{kin} = \frac{1}{2}m\,v^2$ ist proportional zum Quadrat der Geschwindigkeit. Mit $v_1 = \vaO$ und $v_2 = \vbO$ ist
\begin{align}
 \frac{E_2}{E_1} &= \frac{\frac{1}{2}m\,v_2^2}{\frac{1}{2}m\,v_1^2} = \fakF = \frac{\left(\vbO\right)^2}{\left(\vaO\right)^2} = \fak
\end{align}
Die kinetische Energie nimmt also um
\begin{align}
 \frac{E_2-E_1}{E_1} &= \relF = \rel \approx \result{\relP{0}{3}}
\end{align}
zu, sie wird \fakP{0}{3}-mal so gross.
